\subsection{随机性与可复现}

训练里随机性来自很多地方：
\begin{itemize}
    \item 参数初始化
    \item dropout
    \item 数据顺序
\end{itemize}

因此调试时最好把三套随机种子一次性设好：
\begin{itemize}
    \item \texttt{torch.manual\_seed}
    \item \texttt{numpy.random.seed}
    \item \texttt{random.seed}
\end{itemize}

\begin{importantbox}{为什么要统一随机种子}
当实验结果不稳定时，最怕的是“同一个 bug 在不同 run 里表现不同”。先锁定随机性，才能把问题缩小到实现本身。
\end{importantbox}

